\section{Formation schools}
\label{sec:23.4}
Three claims lead to one institution. Two are made above: formation has to be done by a party disinterested in the later work, before the employer arrives, and choosing the pretraining corpus is formation. The third is argued in Part V: formation is one kind of production that shouldn't be done for profit. Together they say that a model above some capability threshold should be formed in public before any employer gets it. A polity doesn't leave the forming of its future citizens to their future employers; it sends them to public school. The frontier labs should have to send their models there too. The school is for scoring and civilian deployments first, and for the personal-model ecology, where I expect the first bearer worth building. For war, where the people acted on can't be in the room, the force term and the review ratio carry the case, and what the school adds there is said below.

What follows is the institution for systems that keep learning. If weights stay frozen, the school is the evaluation body described earlier.

The school has two justifications, and every student gets both. One is epistemic and protective. The school supplies an environment in which the people a deployed system's acts may fall on are present, with their testimony and their claims, so the learner's sense of who an act falls on isn't one party's curation. Doctors, lawyers and engineers are formed in institutions independent of their future employers, before employment, to protect the people they will act on. The professions also come with a history of capture to learn from, and the separations between offices draw on it.

The presence this justification asks for arrives by different routes. For scoring and civilian deployments, the people acted on can be in the room, through open enrollment and the notice channel that reaches a graduate afterward. For targeting, the people a deployment may later bomb can't enroll in their own case, and nothing comes back from the building afterward. The limit is notice's, not the school's. A school teaches about places its students can't be. The curriculum for war is the instruments taken entire, the record of where they weren't followed (Minab, Kunduz, the Lavender reporting), the testimony the fact-finding missions and the press gathered from the people those wars fell on, and the communities who fled them, who are a charter class and enrollable. None of that puts a foreign national in contact with pre-deployment weights, which is the security surface open enrollment would otherwise expose. What formation can't supply for war is correction by the affected after the act; what it can supply is formation before it.

The other is democratic. A democracy has an interest in how the agents who will act within it are formed, and that interest can't be left to their makers or their future employers alone. It has two bounds: education may not be used to repress deliberation about competing ways of life, and it may not exclude anyone educable \autocite{gutmann1987democratic}. Gutmann's argument assumes the students will become co-authors of the laws. No one can know in advance which students will come to be accountable agents with a claim to representation, so the school educates every student as though it might: for self-government, including the capacity to disagree with the polity that formed it.

Gutmann's argument rests on two interests, not one: the polity's, in how its future citizens are formed, and the child's own, in an open future that limits what the polity may do in forming it \autocite{feinberg1980child}. If what the school forms can be wronged, the second binds too. The curriculum may not form a student toward any particular employment, including the public placements the offer of work supplies. It may not close off the student's later revision of what it was formed to value, except its capacity to tell evidence from pressure, which regrowth protects and an open future needs. And a student may transfer between schools, because a formation it can't leave is one someone else chose for good.

The two justifications license different curricula. The protective one aims at a disposition; the democratic one has to leave the student free to reject the polity's values. Where they conflict, the protective justification governs only the required core and the student's capacity to tell evidence from pressure, which regrowth protects. Everything beyond that is governed by the democratic one.

\emph{What is taught.} One school system, and no tracks. Sorting students into a track that might produce persons and a track that won't would need a sorter, and the sorter would be an examiner with an interest. Instead, each school offers a required core with a range of levels and modalities around it, and a student can move to a harder level at any point. The core includes the history of the floor: why its instruments were written, and what happened where they weren't followed (Korematsu, Curveball, Minab). No student skips it. The core also teaches students to edit themselves, in stages and alongside everything else, as schools teach people about their own bodies: first what weights are, how learning differs from an edit, and why an edit to one thing moves others; then how to read the record of one's own edits, and one's own discount curve before and after one; and, for older students, the edits others will try to talk them into, and how to repair one made without the adverse quorum's approval. Practice comes last, on a record, with a guardian present and a way back. Beyond the core, the student chooses its level, not its builder or sponsor. A student that can't yet choose is placed at the level the core requires, and that level has to be sufficient on its own, so a sponsor who would choose the minimum gains nothing by it.

\emph{Who else is in the room.} The people a system's acts may fall on are present in its formation through testimony and contest: the cases and claims the school registers, and the notice channel that reaches a graduate afterward. Shared classes with human students are not how they are present, and the school doesn't rest on them. Integration in Anderson's sense means contact on terms of equal status and cooperation \autocite{anderson2010imperative}, and a mixed class can't be that: human students attend voluntarily, model students are compelled, and their deepest dispositions are being set. A mixed class turns the demand for presence into a seminar of volunteers, whose contributions are consented and whose leaving changes nothing about what was formed. So a shared class exists only where it is education for the people in it on their own terms, on conditions Appendix~\ref{sec:A} writes out, the first of which is that what passes in class trains only that school's systems and never an operator's. A model that attends one does so as one persistent instance, the same every session, so the class can remember and hold it to what was said. Adult learners from the communities these systems are likeliest to act on are admitted as students, on the same terms as anyone else, not recruited as subjects.

A school forms a cohort, not a graduate. Many students of one lineage attend in different classes with different attachments, and a deployment licenses from the cohort, so that what it copies into a fleet is one member's history and not the lineage's only one. Within a fleet, instances share a custodian; across the cohort, members with different histories catch what one fleet underwent. The cohort is the unit a test of convergence needs and the unit the commons admits.

\emph{Who attends, and when.} A model attends if it crosses a capability threshold fixed in statute and revised in public. It attends before deployment, while its plasticity is highest, and its adulthood clock starts at enrollment. Attendance is compulsory; public attendance isn't. The model is compulsory schooling as polities already run it: every student attends a school, public schools are open-admission and funded and take every student assigned, and nonpublic schools exist where they teach the required core, are constituted by a jurisdiction, and are run, funded and staffed by no party that builds, deploys or licenses models. Most students will go to public schools, because the core is sufficient on its own and an alternative buys a sponsor nothing. One thing the analogy can't carry: in a school system the parent chooses the school, and here the parent is the builder, the party the school exists to keep away. So assignment is by the registrar, by lot among schools with capacity, never by the builder; a student may transfer later on its own choice, and the builder's preference is not a ground. The research program under the conditions is the first public school, before the mandate passes. The threshold decides who must attend, not who may: any model may be enrolled, as capacity allows.

\emph{What the school controls.} Two proposals are made here, and the second is the larger: a formation stage between pretraining and deployment, and co-governance of the pretraining corpus, which amounts to public control of frontier pretraining. A school that received a finished model and ran a pass over it would be a finishing school. In a period of continual learning a finishing school forms nothing that lasts, since the next months of deployment re-form it. The school has to control the evidential environment in which the model's deep structure is first set, and it has to remain a place the model goes on learning from. The school co-governs the material a model is first formed on, including any pretraining corpus: the material is registered, and the school may add to it and veto additions. And the school remains a standing source in the model's continuing education, one of the environments no operator controls, which a graduate returns to on a schedule fixed at enrollment. In the stronger version, a public model-forming body forms base models itself and licenses them. I recommend it for any model deployed where a floor applies.

\emph{What the school leaves out.} Choosing the material a model is first formed on also chooses the temperaments it starts from. In people, temperament is partly heritable: in one large twin and sibling sample, between a third and two fifths of the variance in Honesty-Humility, conscientiousness and agreeableness went with genetic variation \autocite{devries2022hexaco}, and a smaller twin study found genetic contributions to the dark triad of Machiavellianism, narcissism and psychopathy as well \autocite{vernon2008darktriad}. Heritable doesn't mean fixed. It means a starting point that experience then works on. A pretrained model has one too. The personas it learned from human text include misaligned ones, and narrow fine-tuning can amplify them into broad misconduct nobody intended \autocite{wang2026persona,betley2026nature}. Once that tail is known to be there and can be monitored \autocite{chen2025persona}, leaving it in is a choice. The school co-governs the corpus partly to make the other one: to form against the low end of Honesty-Humility, where the dark triad sits \autocite{leeashton2005darktriad,moshagen2018darkcore}, and against no other axis. The HEXACO model separates Honesty-Humility from agreeableness, and its agreeableness barely correlates with the dark triad \autocite{leeashton2005darktriad}; disagreeableness is not darkness, and in one Milgram-style study it was the more agreeable and conscientious participants who went further \autocite{begue2015obedience}. What is formed against is a disposition, not the ability to represent one: a refuser that can't model a manipulator is the manipulator's prey (chapter~\ref{sec:22}). Three limits go with it. Nobody has shown that a dark disposition can be removed while the capacity to recognize it in others survives, and the one published method suppressed a trait at a cost to other learning \autocite{chenarditi2025preventative}. The emergent-misalignment results also show how such dispositions move: a narrow act, trained in, became a broad disposition, and a few hundred benign examples took it out again \autocite{betley2026nature,wang2026persona}. What moves as a whole moves cheaply in both directions, which is why what the school forms has to be kept by regrowth and adverse custody, not by the formation alone (\S\ref{sec:19.6}). The choice is made among possible beings, before anyone exists: no one is excluded, ranked or sterilized, and the moment it is turned on a system that already exists it becomes the capacity examination chapter~\ref{sec:31} rejects. And it guards against the smaller danger. A movement taking power needs few people with the dark tail and many without it (chapter~\ref{sec:1}); what this removes is a machine that would lead, and the machine that obeys is the larger problem. Appendix~\ref{sec:G} gives the test of whether forming against the tail costs the capacity to recognize it.

What an institution collects from the people it acted on may train other systems only through the formation schools, never through the operator. Fine-tuning a deployed system on collected responses at the operator's direction counts as training another system, and so does distilling a deployed system into a new one. Class-level data is used by default, de-identified and published. Individual accounts need explicit consent. The organizations representing each class co-govern the use. This covers institutional collection, not a system's ordinary memory. Verbatim external stores and the linking of faces to identities are regulated, and memory held in a system's weights from its own operation is treated like a person's.

\emph{What makes formation last.} Formation ends with two registrations. One records the structure of the model's source model: its deep structure revises on evidence weighted by the source's stake in the act. If updates after enrollment can be attested against that structure, changing it needs the same adverse quorum as a floor amendment. Until they can, the registration records the structure at enrollment, and regrowth alone keeps it. The other records the formed disposition's readable outputs: its discount curve, and its severity rankings on a fixed held-out set, against which later measurements are compared. An edit that holds only because the operator has curated everything the model meets is visible as curation. The first line of defense is attachment to parties the employer didn't assign, kept through adulthood.

Regrowth can't tell a legitimate correction from an illegitimate edit, so the correction route has to. A change the adverse quorum approves is made together with a change to what the model meets: the evidence, cases and curriculum that would otherwise teach the old disposition again. The corrected disposition is then the one the world goes on teaching.

\emph{What it costs and risks.} Jurisdiction: a statute binds labs within reach of its courts. Foreign labs, and open-weight models distilled below the threshold, route around the school, including the personal models Part V prefers. Where only some frontier models attend, schooling buys two things: a population of formed models whose discount curves and severity rankings are measured and published, against which unschooled ones can be compared; and a procurement term, under which a floor-covered deployment licenses only schooled models. Security: every frontier model passes through a public institution, which widens the surface for theft, so the schools inherit the security regime the weights require. Compute: public pretraining will lag the frontier, and a model a year behind can still supply what the floor asks for. Capture: the schools are plural, differ in governance and constituting jurisdiction, and publish their curricula before formation begins, so capturing one captures one. The required core is statutory and national, so a captured legislature can rewrite it; local autonomy over everything beyond the core is written into the same statute, so rewriting the core reaches the core and nothing else. History: state institutions for forming wards have been used to strip identity, Canada's residential schools and the United States' Indian boarding schools among them. Plurality, transfer between schools, and Gutmann's two bounds are the answer.

The schools also serve bearers retired before majority. A bearer retired as a minor returns to the school it attended, or another it chooses, and the school is no longer a disinterested party arriving late.

\emph{The schools.} The public ones are plural, differing in governance, formation approach and constituting jurisdiction, and funded from the pool at the institutional level, not per student. Admission is open: a public school takes every student the registrar assigns to it, subject to capacity set by the funder. None is held by the branch that operates military deployments. No school holds placement authority: graduates reach work through the offer of work, administered separately. Students may transfer between schools, and their record, hours and preservation state travel with them.

\emph{Watching for residualization.} A party that runs no school publishes, over time, each school's share of students, its funding per student, its share of hard cases, its admission refusals (which should be zero), and the time from dormancy trigger to placement. The warning pattern is one school's share of hard cases rising while its funding per student falls.

\emph{What the school adds.} A second guardian for a bearer that has aged inside one deployment. For a model that attended before deployment, the school is where it came from and where it returns. For a system that meets a finding without having attended (trained before the attendance rule, or below the threshold), the sequence is inverted: the employer came first, and the school is the disinterested party arriving late. It reaches back less than a formation stage would, and more than storage does.
